\documentclass[10pt,a4paper]{amsart}
\usepackage{amsmath,amssymb,amsthm,graphicx}
% Retained minimal preamble for the Stacks Project chapter flat.tex (More on Flatness).
% Verbatim excerpt of the chapter preamble: Xy-pic setup (native preamble lines 14-18) and the native
% theorem declarations with the mathematical macros (native preamble lines 154-217). Copied mechanically;
% nothing added, removed or redefined. The wrapper class and the amsmath/amssymb/amsthm/graphicx base
% packages are supplied by the assembler, so the class/externaldocument part of the native preamble is omitted.
% For commutative diagrams we use Xy-pic
\usepackage[all]{xy}

% We use 2cell for 2-commutative diagrams.
\xyoption{2cell}
% Theorem environments.
%
\theoremstyle{plain}
\newtheorem{theorem}[subsection]{Theorem}
\newtheorem{proposition}[subsection]{Proposition}
\newtheorem{lemma}[subsection]{Lemma}

\theoremstyle{definition}
\newtheorem{definition}[subsection]{Definition}
\newtheorem{example}[subsection]{Example}
\newtheorem{exercise}[subsection]{Exercise}
\newtheorem{situation}[subsection]{Situation}

\theoremstyle{remark}
\newtheorem{remark}[subsection]{Remark}
\newtheorem{remarks}[subsection]{Remarks}

\numberwithin{equation}{subsection}

% Macros
%
\def\lim{\mathop{\mathrm{lim}}\nolimits}
\def\colim{\mathop{\mathrm{colim}}\nolimits}
\def\Spec{\mathop{\mathrm{Spec}}}
\def\Hom{\mathop{\mathrm{Hom}}\nolimits}
\def\Ext{\mathop{\mathrm{Ext}}\nolimits}
\def\SheafHom{\mathop{\mathcal{H}\!\mathit{om}}\nolimits}
\def\SheafExt{\mathop{\mathcal{E}\!\mathit{xt}}\nolimits}
\def\Sch{\mathit{Sch}}
\def\Mor{\mathop{\mathrm{Mor}}\nolimits}
\def\Ob{\mathop{\mathrm{Ob}}\nolimits}
\def\Sh{\mathop{\mathit{Sh}}\nolimits}
\def\NL{\mathop{N\!L}\nolimits}
\def\CH{\mathop{\mathrm{CH}}\nolimits}
\def\proetale{{pro\text{-}\acute{e}tale}}
\def\etale{{\acute{e}tale}}
\def\QCoh{\mathit{QCoh}}
\def\Ker{\mathop{\mathrm{Ker}}}
\def\Im{\mathop{\mathrm{Im}}}
\def\Coker{\mathop{\mathrm{Coker}}}
\def\Coim{\mathop{\mathrm{Coim}}}

% Boxtimes
%
\DeclareMathSymbol{\boxtimes}{\mathbin}{AMSa}{"02}

%
% Macros for moduli stacks/spaces
%
\def\QCohstack{\mathcal{QC}\!\mathit{oh}}
\def\Cohstack{\mathcal{C}\!\mathit{oh}}
\def\Spacesstack{\mathcal{S}\!\mathit{paces}}
\def\Quotfunctor{\mathrm{Quot}}
\def\Hilbfunctor{\mathrm{Hilb}}
\def\Curvesstack{\mathcal{C}\!\mathit{urves}}
\def\Polarizedstack{\mathcal{P}\!\mathit{olarized}}
\def\Complexesstack{\mathcal{C}\!\mathit{omplexes}}
% \Pic is the operator that assigns to X its picard group, usage \Pic(X)
% \Picardstack_{X/B} denotes the Picard stack of X over B
% \Picardfunctor_{X/B} denotes the Picard functor of X over B
\def\Pic{\mathop{\mathrm{Pic}}\nolimits}
\def\Picardstack{\mathcal{P}\!\mathit{ic}}
\def\Picardfunctor{\mathrm{Pic}}
\def\Deformationcategory{\mathcal{D}\!\mathit{ef}}

% --- appended 2026-10-12 (Chief): the project own definitions of the reference/slogan/history
% display environments, copied mechanically from the shared Stacks preamble lines 8-12
% (\usepackage{verbatim} plus the three \newenvironment lines, which the chapter preamble keeps
% verbatim from the project preamble). Nothing added, removed or redefined: these environments
% carry no mathematical content and the project itself typesets them as comments. Without these
% three lines a counted statement whose own text uses them fails with Environment slogan undefined.
\usepackage{verbatim}
\newenvironment{reference}{\comment}{\endcomment}
\newenvironment{slogan}{\comment}{\endcomment}
\newenvironment{history}{\comment}{\endcomment}
\usepackage{cleveref}
\crefname{originallemma}{original lemma}{original lemmas}
\Crefname{originallemma}{Original lemma}{Original lemmas}
\makeatletter
\newlabel{algebra-lemma-flat-localization}{{00HT}{source}{}{sourcepointer-1}{}}
\newlabel{algebra-lemma-flat-localization@cref}{{[originallemma][0][]00HT}{[1][0][]source}{}{}{}}
\makeatother
\crefname{originallemma}{original lemma}{original lemmas}
\Crefname{originallemma}{Original lemma}{Original lemmas}
\makeatletter
\newlabel{algebra-lemma-integral-going-up}{{00GU}{source}{}{sourcepointer-2}{}}
\newlabel{algebra-lemma-integral-going-up@cref}{{[originallemma][0][]00GU}{[1][0][]source}{}{}{}}
\makeatother
\crefname{originalreference}{original reference}{original references}
\Crefname{originalreference}{Original reference}{Original references}
\makeatletter
\newlabel{situation-flat-at-point}{{05MK}{source}{}{sourcepointer-3}{}}
\newlabel{situation-flat-at-point@cref}{{[originalreference][0][]05MK}{[1][0][]source}{}{}{}}
\makeatother
\title{The flat at a point functor is unchanged by a finite local extension carrying an isomorphic module}
\author{The Stacks Project Authors}
\date{stacks-project:flat.tex:local-snapshot-20261011; GNU FDL 1.2 or later}
\begin{document}
\maketitle
\noindent\textit{Editorial note.} All delivered mathematical text is the author's own, copied from the pinned Stacks Project chapter file flat.tex; no statement or proof step was rewritten, reordered or removed.\par\medskip
\noindent\textit{Editorial note.} Preamble provenance: the chapter's own retained minimal preamble (sources/stacks-flat/retained-preamble.tex, lines 1-85) is a verbatim excerpt of the native Xy-pic setup and of the native theorem declarations and mathematical macros, followed by the project's own reference/slogan/history display-environment definitions copied mechanically from the shared preamble; the wrapper class and the amsmath/amssymb/amsthm/graphicx base packages are supplied by the assembler, and extra\_packages is left empty because everything needed is already in the retained preamble. No macro was invented or redefined.\par\medskip
\noindent\textit{Editorial note.} Same-chapter pointers are located at their \textbackslash{}label lines inside the pinned chapter file; cross-chapter pointers are located in the project's own tag list staged in the snapshot (sources/stacks-chow/tags.txt), at the verbatim lines carrying their canonical tags.\par\medskip
\noindent\textit{Allowed standard prerequisites (editorial).}\begin{itemize}
\item The situation defining the flat at a point functor. Same chapter flat.tex, situation labelled situation-flat-at-point (statement, proof lines 6069 and 6079). Fixes the category over which the functors are defined and the notation for the base changes of the algebras and modules.
\item Flatness of a module over a local ring is detected at maximal ideals of the localisation. Algebra, label algebra-lemma-flat-localization, invoked at proof line 6097. Allows the flatness of the localised module to be checked after localising further at maximal ideals, which is where the finite extension lets the maximal ideals be compared.
\item The maximal ideals of a finite local extension correspond to the primes over the given prime. Algebra, label algebra-lemma-integral-going-up, invoked at proof line 6101. Identifies the maximal ideals of the localised extension with the primes lying over the given prime, so that the vanishing locus computed upstairs matches the image of the one downstairs.
\end{itemize}\medskip
\ifdefined\hypertarget\hypertarget{sourcepointer-1}{}\fi
\noindent\textit{Source reference.} Original lemma 00HT: sources/stacks-chow/tags.txt, source lines 655--655. Navigation pointer for \textbackslash{}ref\{algebra-lemma-flat-localization\} used inside the counted span (line 6097); the argument that consumes it is retained in the span.\par\medskip
\ifdefined\hypertarget\hypertarget{sourcepointer-2}{}\fi
\noindent\textit{Source reference.} Original lemma 00GU: sources/stacks-chow/tags.txt, source lines 620--620. Navigation pointer for \textbackslash{}ref\{algebra-lemma-integral-going-up\} used inside the counted span (line 6101); the argument that consumes it is retained in the span.\par\medskip
\ifdefined\hypertarget\hypertarget{sourcepointer-3}{}\fi
\noindent\textit{Source reference.} Original reference 05MK: sources/stacks-flat/pinned-original-source.tex, source lines 6008--6008. Navigation pointer for \textbackslash{}ref\{situation-flat-at-point\} used inside the counted span (line 6069); the argument that consumes it is retained in the span.\par\medskip
\begin{lemma}
\label{lemma-flat-at-point-finite}
In Situation \ref{situation-flat-at-point}. Let $B \to C$ is a local map of
local $A$-algebras and $N$ a $C$-module. Denote
$F'_{lf} : \mathcal{C} \to \textit{Sets}$ the functor associated to the pair
$(C, N)$. If $M \cong N$ as $B$-modules and $B \to C$ is finite, then
$F_{lf} = F'_{lf}$.
\end{lemma}

\begin{proof}
Let $A'$ be an object of $\mathcal{C}$. Set $C' = C \otimes_A A'$
and $N' = N \otimes_A A'$ similarly to the definitions of $B'$, $M'$ in
Situation \ref{situation-flat-at-point}.
Note that $M' \cong N'$ as $B'$-modules.
The assumption that $B \to C$ is finite has two consequences:
(a) $\mathfrak m_C = \sqrt{\mathfrak m_B C}$ and (b)
$B' \to C'$ is finite. Consequence (a) implies that
$$
V(\mathfrak m_{A'}C' + \mathfrak m_C C')
=
\left(
\Spec(C') \to \Spec(B')
\right)^{-1}V(\mathfrak m_{A'}B' + \mathfrak m_B B').
$$
Suppose $\mathfrak q \subset V(\mathfrak m_{A'}B' + \mathfrak m_B B')$.
Then $M'_{\mathfrak q}$ is flat over $A'$ if and only if
the $C'_{\mathfrak q}$-module $N'_{\mathfrak q}$ is flat over $A'$
(because these are isomorphic as $A'$-modules) if and only if
for every maximal ideal $\mathfrak r$ of $C'_{\mathfrak q}$
the module $N'_{\mathfrak r}$ is flat over $A'$ (see
Algebra, Lemma \ref{algebra-lemma-flat-localization}).
As $B'_{\mathfrak q} \to C'_{\mathfrak q}$ is finite by (b),
the maximal ideals of $C'_{\mathfrak q}$ correspond exactly
to the primes of $C'$ lying over $\mathfrak q$ (see
Algebra, Lemma \ref{algebra-lemma-integral-going-up})
and these primes are all contained in
$V(\mathfrak m_{A'}C' + \mathfrak m_C C')$ by
the displayed equation above. Thus the result of the lemma holds.
\end{proof}

\end{document}
